\chapter{پژوهش‌های مرتبط و روش‌های مبنا}

\section{مقدمه و چارچوب طبقه‌بندی}

هدف این فصل، مرور پژوهش‌های مرتبط با تشخیص توضیح‌پذیر سلامت روان و بررسی روش‌هایی است که مبنای بازپیاده‌سازی و آزمایش‌های این پژوهش قرار می‌گیرند. ابتدا کارهای مرتبط با وظیفه، از طبقه‌بندی متن تا تولید توضیح و استفاده از دانش دامنه‌ای، مرور می‌شوند. سپس روش‌های بازیابی افزوده و گراف‌محور از نظر مسئله اولیه، نوع شواهد و نحوه تولید پاسخ بررسی می‌شوند تا شرایط لازم برای انطباق آن‌ها با سلامت روان روشن گردد.

روش‌های منتخب بر اساس محل دخالت گراف در فرایند پاسخ‌گویی به سه خانواده تقسیم می‌شوند. در خانواده بازیابی‌محور، بخشی از گراف یا اسناد مرتبط پیش از تولید پاسخ انتخاب و به مدل داده می‌شود؛ \lr{G-Retriever}، \lr{HippoRAG 2} و در مقیاسی محدودتر \lr{LinearRAG} در این دسته بررسی می‌شوند. در خانواده عامل‌محور، مدل زبانی در چند مرحله با گراف تعامل می‌کند و مسیر جست‌وجو را شکل می‌دهد؛ \lr{ToG} و \lr{ToG-2} نمایندگان این الگو هستند. در نهایت، روش \lr{GCR} به عنوان یک الگوی قیدگذارانه بررسی می‌شود که ساختار گراف را در فرایند تولید مسیر استدلال وارد می‌کند.

این طبقه‌بندی به معنای برتری ذاتی یک خانواده بر خانواده دیگر نیست. هر الگو برای هدف و تنظیم آزمایشی متفاوتی طراحی شده و میان پوشش شواهد، انعطاف استدلال، وفاداری مسیر و هزینه اجرا تعادل خاصی برقرار می‌کند. از این رو، مقایسه روش‌ها علاوه بر کیفیت پاسخ، به نیاز آن‌ها به آموزش، تعداد تعامل‌ها با مدل زبانی، نحوه انتخاب زیرگراف یا متن مرتبط، امکان ردیابی شواهد، زمان پاسخ و هزینه متنی توجه خواهد داشت. این معیارها مبنای تنظیم مجموعه ارزیابی مشترک در فصل‌های بعد هستند.

پس از مرور پژوهش‌های سلامت روان، بازیابی افزوده به عنوان راهی برای دسترسی مدل به دانش بیرونی معرفی می‌شود و گذار از بازیابی متنی به گراف‌محور توضیح داده می‌شود. پیش از بررسی شیوه‌های استفاده از گراف، روش \lr{KGGen} برای ساخت گراف دانش از متن معرفی می‌شود. سپس روش‌های بازیابی‌محور و عامل‌محور معرفی و مقایسه می‌شوند و \lr{GCR} از نظر قیدگذاری ساختاری بررسی می‌گردد. بخش پایانی نیز نحوه انتقال این روش‌ها به وظایف منتخب سلامت روان و معیارهای ارزیابی مشترک را جمع‌بندی می‌کند.

\section{پژوهش‌های تشخیص توضیح‌پذیر سلامت روان}

پژوهش‌های این حوزه از نظر نوع خروجی و نحوه استفاده از دانش متفاوت‌اند. برخی، متن یا سابقه فعالیت کاربر را به یک برچسب نگاشت می‌کنند؛ برخی همراه با برچسب، توضیح زبانی تولید می‌کنند؛ و گروهی نیز معیارهای تشخیصی یا منابع دانش را در پردازش وارد می‌کنند. مرور این کارها مشخص می‌کند که روش‌های گرافی در این پایان‌نامه قرار است در چه وظیفه‌ای و با چه خروجی‌هایی ارزیابی شوند.

\subsection{از روش‌های پیش‌بینی تا مدل‌های زبانی بزرگ}

در کارهای اولیه، تشخیص افسردگی عمدتاً به صورت یک مسئله طبقه‌بندی از داده‌های شبکه‌های اجتماعی صورت‌بندی شده است. شن و همکاران از یادگیری فرهنگ‌نامه چندوجهی برای ترکیب اطلاعات کاربران استفاده کرده‌اند و اورابی و همکاران، یادگیری عمیق را برای تشخیص افسردگی کاربران توییتر بررسی کرده‌اند \cite{shen2017depression,orabi2018deep}. پژوهش \lr{SenseMood} نیز در همین مسیر به تشخیص افسردگی در شبکه‌های اجتماعی می‌پردازد \cite{lin2020sensemood}. این مطالعات مبنایی برای یادگیری الگوهای مرتبط با برچسب فراهم می‌کنند، در حالی که تولید توضیح زبانی برای هر تصمیم، خروجی اصلی آن‌ها نیست.

روش \lr{AMMNet} مسئله پیش‌بینی افسردگی و اضطراب در \lr{Reddit} را با یادگیری چندوظیفه‌ای دنبال می‌کند \cite{sarkar2023predicting}. در این روش، نمایش‌های متنی مشترک با ویژگی‌های موضوعی اختصاصی وظایف ترکیب می‌شوند. استفاده از اطلاعات مشترک، امکان بهره‌گیری از ارتباط میان وظایف را فراهم می‌کند و ویژگی‌های موضوعی نیز به بررسی محتوای مؤثر در پیش‌بینی کمک می‌کنند. این نوع تفسیر ویژگی‌ها با توضیح آزاد زبانی تفاوت دارد، اما برای مقایسه شیوه‌های ارائه شواهد تصمیم مفید است.

با ورود مدل‌های زبانی بزرگ، امکان تولید هم‌زمان برچسب و توضیح مطرح شد. یانگ و همکاران در پژوهش «به سوی تحلیل تفسیرپذیر سلامت روان با مدل‌های زبانی بزرگ»، تنظیم‌های بدون نمونه، چندنمونه‌ای و استدلال زنجیره‌ای را بررسی کرده‌اند \cite{yang2023towards}. مدل در این تنظیم‌ها با استفاده از دستور و نمونه‌های داخل ورودی به وظیفه هدایت می‌شود و برای هر وظیفه الزاماً آموزش مجدد نمی‌بیند. ارزیابی انسانی توضیح‌ها نیز بخشی از مطالعه است؛ بنابراین، این کار علاوه بر پیش‌بینی برچسب، کیفیت توضیح را به عنوان خروجی قابل سنجش مطرح می‌کند. چنین تنظیم‌هایی می‌توانند مبنای مقایسه با روش‌هایی باشند که دانش بیرونی دریافت می‌کنند.

\subsection{تنظیم دامنه‌ای و استفاده از دانش برای تولید توضیح}

روش \lr{MentalLLaMA} تحلیل توضیح‌پذیر سلامت روان را از طریق تنظیم دستوری مدل دنبال می‌کند \cite{yang2024mentallama}. داده‌های چند وظیفه در مجموعه \lr{IMHI} به قالب دستور و پاسخ تبدیل شده‌اند و پاسخ‌ها شامل برچسب و توضیح‌اند. توضیح‌ها با کمک مدل زبانی تولید و با ارزیابی‌های خودکار و انسانی بررسی شده‌اند. آموزش بر این داده‌ها، مدل را با صورت وظایف و نحوه پاسخ‌گویی آشنا می‌کند. ارتباط مستقیم این کار با پایان‌نامه حاضر، استفاده از \lr{IMHI} و ارزیابی جداگانه پیش‌بینی و توضیح است؛ در اینجا نقش افزودن شواهد گرافی به همین نوع خروجی بررسی خواهد شد.

چارچوب \lr{DORIS} بر تشخیص افسردگی از پست‌های کاربران تمرکز دارد و دانش تشخیصی را در استخراج ویژگی‌ها وارد می‌کند \cite{lan2025doris}. مدل زبانی، متن‌ها را از نظر نشانه‌های مرتبط با معیارهای تشخیصی بررسی و روند زمانی حالات خلقی را خلاصه می‌کند. این اطلاعات همراه با نمایش محتوای پست‌ها برای آموزش یک طبقه‌بند به کار می‌روند. توضیح پیش‌بینی نیز بر پایه نشانه‌ها و تحلیل روند خلقی تولید می‌شود. در نتیجه، مدل زبانی هم در آماده‌سازی شواهد و هم در بیان توضیح نقش دارد و تصمیم نهایی از یک مؤلفه طبقه‌بندی حاصل می‌شود. این پژوهش نمونه‌ای از استفاده صریح از دانش دامنه‌ای است، هرچند ورودی آن سابقه پست‌های کاربر است و تنظیم آن با تمام وظایف \lr{IMHI} یکسان نیست.

در سطح منابع دانش، گائو و همکاران گراف \lr{MDKG} را با استفاده از متون علمی، پایگاه‌های زیست‌پزشکی و مدل‌های زبانی ساخته‌اند \cite{gao2025mdkg}. این کار اطلاعات مرتبط با اختلالات روانی را در یک ساختار مشترک سازمان‌دهی می‌کند و کاربرد نمایش‌های حاصل را در پیش‌بینی نیز می‌سنجد. نقش آن در مرور حاضر، معرفی یک مسیر برای فراهم کردن دانش دامنه‌ای است؛ انتخاب و آماده‌سازی منبع گراف برای آزمایش‌های این پایان‌نامه در فصل روش تحقیق مشخص خواهد شد.

این پژوهش‌ها سه راه برای پشتیبانی از تصمیم را نشان می‌دهند: یادگیری از داده‌های برچسب‌دار، آموزش نحوه تولید توضیح و استفاده از دانش دامنه‌ای. بازیابی افزوده راه دیگری برای دسترسی به دانش فراهم می‌کند که در آن، شواهد متناسب با هر ورودی انتخاب و به مدل داده می‌شوند. بخش بعد این سازوکار را توضیح می‌دهد و سپس به نقش روابط گرافی در انتخاب و استفاده از شواهد می‌پردازد.

\section{از بازیابی افزوده متنی تا بازیابی افزوده گراف‌محور}

در الگوی متداول بازیابی افزوده، منبع بیرونی مجموعه‌ای از اسناد متنی است که پس از قطعه‌بندی و تبدیل به نمایش‌های برداری، نمایه‌سازی می‌شوند. برای هر پرسش، چند قطعه با بیشترین شباهت معنایی بازیابی و همراه با پرسش به مدل زبانی داده می‌شوند. این الگو امکان استفاده از اطلاعات بیرون از پارامترهای مدل را بدون آموزش دوباره فراهم می‌کند و به دلیل سادگی و قابلیت اجرا روی مجموعه‌های بزرگ، مبنای بسیاری از سامانه‌های بازیابی افزوده قرار گرفته است \cite{lewis2020retrieval}.

بازیابی متنی بر دو فرض اصلی استوار است: نخست آن‌که قطعه مرتبط را بتوان با شباهت میان پرسش و متن تشخیص داد و دوم آن‌که اطلاعات لازم برای پاسخ در همان قطعه یا در تعداد محدودی از قطعه‌های برتر وجود داشته باشد. این فرض‌ها برای پرسش‌های مستقیم و واقعیت‌محور مناسب‌اند. افزون بر این، اسناد می‌توانند مستقل از یکدیگر به نمایه افزوده یا از آن حذف شوند و هزینه جست‌وجو نیز با استفاده از نمایه‌های برداری قابل کنترل است.

در پرسش‌های چندگامی، شواهد ممکن است میان چند قطعه یا سند توزیع شده باشند و ارتباط آن‌ها فقط از طریق یک یا چند مفهوم واسط آشکار شود. در این حالت، قطعه‌ای که به‌تنهایی شباهت زیادی با پرسش ندارد ممکن است یکی از حلقه‌های لازم برای رسیدن به پاسخ باشد. انتخاب مستقل قطعه‌ها چنین ارتباطی را مستقیماً نمایش نمی‌دهد و ترکیب آن‌ها به مدل زبانی واگذار می‌شود. همین مسئله سبب شده است برخی پژوهش‌ها ساختارهایی را به بازیابی اضافه کنند که پیوند میان بخش‌های مختلف دانش را نیز در مرحله جست‌وجو لحاظ می‌کنند \cite{ma2025tog2,gutierrez2025hipporag2}.

در بازیابی افزوده گراف‌محور، واحدهای اطلاعاتی به صورت گره و ارتباط میان آن‌ها به صورت یال نمایش داده می‌شوند. گره می‌تواند یک موجودیت، گزاره، جمله یا قطعه متن باشد و یال نیز رابطه دانشی، هم‌رخدادی یا پیوند معنایی میان دو گره را نشان دهد. بنابراین، خروجی بازیاب لزوماً مجموعه‌ای از قطعه‌های مستقل نیست و می‌تواند شامل چند سه‌تایی، یک مسیر یا یک زیرگراف متصل باشد. برای نمونه، \lr{G-Retriever} بخشی از گراف متنی را با توجه به پرسش انتخاب می‌کند و ساختار همان بخش را همراه با محتوای متنی در اختیار مدل قرار می‌دهد \cite{he2024gretriever}.

گذار به بازیابی گراف‌محور به معنای جایگزینی کامل متن با گراف نیست. گراف روابط و مسیرهای ارتباطی را به‌صورت صریح نگهداری می‌کند، در حالی که متن جزئیات و زمینه‌ای را در بر دارد که ممکن است در سه‌تایی‌های گراف ثبت نشده باشد. به همین دلیل، برخی روش‌ها از گراف برای هدایت بازیابی متن استفاده می‌کنند و برخی دیگر، شواهد گرافی و متنی را به‌صورت ترکیبی به مدل می‌دهند. \lr{ToG-2} نمونه‌ای از ادغام رفت‌وبرگشتی این دو منبع است که در آن نتایج پیمایش گراف، بازیابی متن را هدایت می‌کند و متن بازیابی‌شده نیز بر ادامه پیمایش اثر می‌گذارد \cite{ma2025tog2}.

در نتیجه، بازیابی متنی و گراف‌محور در نحوه سازمان‌دهی دانش، واحد بازیابی و میزان استفاده از روابط ساختاری تفاوت دارند. برای استفاده از روابط، ابتدا باید مشخص شود گراف چگونه از منابع متنی ساخته می‌شود و صورت‌های مختلف یک مفهوم چگونه در آن یکسان می‌شوند. از این رو، بخش بعد روش \lr{KGGen} را برای ساخت گراف دانش از متن معرفی می‌کند. پس از آن، شیوه‌های انتخاب شواهد و پیمایش گراف در روش‌های بازیابی‌محور و عامل‌محور بررسی خواهند شد.

\section{ساخت گراف دانش از متن: روش \lr{KGGen}}

ساخت گراف از متن، مرحله آماده‌سازی دانش برای بازیابی و استدلال است. در این مرحله باید مفاهیم و روابط متن استخراج شوند و اطلاعات اسناد مختلف در یک ساختار قابل جست‌وجو قرار گیرند. برای مثال، ثبت یک مفهوم با چند نام مستقل می‌تواند ارتباط شواهد پراکنده را مخفی کند. در مقابل، ادغام دو مفهوم متمایز ممکن است رابطه‌ای گمراه‌کننده ایجاد کند. بنابراین، کیفیت گراف تنها به تعداد سه‌تایی‌های استخراج‌شده وابسته نیست و نحوه یکسان‌سازی آن‌ها نیز اهمیت دارد.

روش \lr{KGGen} سه مرحله دارد: استخراج، تجمیع و یکسان‌سازی موجودیت‌ها و روابط \cite{mo2025kggen}. ابتدا مدل زبانی فهرست موجودیت‌های هر متن را استخراج می‌کند و سپس با دریافت متن و این فهرست، سه‌تایی‌ها را می‌سازد. در مرحله تجمیع، گراف‌های حاصل از منابع مختلف ترکیب و موارد تکراری یکسان حذف می‌شوند. این مرحله به فراخوانی مدل زبانی نیاز ندارد.

در مرحله سوم، نام‌های متفاوتی که به یک موجودیت یا رابطه اشاره دارند شناسایی می‌شوند. در نسخه دوم مقاله، نمایش‌های برداری برای خوشه‌بندی به کار می‌روند و نامزدهای مشابه با ترکیب جست‌وجوی واژگانی و معنایی پیدا می‌شوند. مدل زبانی هم‌ارزی نامزدها را بررسی و برای موارد تکراری یک نام نماینده انتخاب می‌کند. این فرایند جداگانه بر موجودیت‌ها و روابط اعمال می‌شود تا تنوع نام‌گذاری، گراف را بی‌جهت پراکنده نکند؛ صرف شباهت معنایی نیز نباید به ادغام مفاهیم متفاوت منجر شود.

مقاله برای ارزیابی استخراج، مجموعه ارزیابی \lr{MINE} را معرفی می‌کند که حفظ اطلاعات متن و کاربرد گراف در بازیابی افزوده را می‌سنجد. در این پایان‌نامه، \lr{KGGen} به عنوان یک روش ساخت گراف بررسی می‌شود؛ انتخاب منابع سلامت روان و کنترل کیفیت گراف حاصل، بخشی از طراحی آزمایش خواهد بود. پس از آماده‌سازی چنین ساختاری، مسئله بعدی انتخاب بخش مرتبط آن برای هر ورودی است که در بخش بعد بررسی می‌شود.

\section{روش‌های گرافی بازیابی‌محور: سازوکار و ویژگی‌ها}

در روش‌های گرافی بازیابی‌محور، انتخاب شواهد پیش از تولید پاسخ و به وسیله یک مؤلفه بازیاب انجام می‌شود. مدل زبانی زیرگراف یا قطعه‌های متنی حاصل از بازیابی را دریافت می‌کند، اما مستقیماً درباره گام بعدی پیمایش تصمیم نمی‌گیرد. روش‌های این خانواده در سه مؤلفه با یکدیگر تفاوت دارند: گراف از پیش موجود است یا از متن ساخته می‌شود، امتیاز ارتباط پرسش با گره‌ها و یال‌ها چگونه محاسبه می‌گردد، و خروجی بازیاب به صورت زیرگراف یا مجموعه‌ای از قطعه‌های متن در اختیار مدل قرار می‌گیرد.

روش \lr{G-Retriever} برای پرسش و پاسخ روی گراف‌های متنی طراحی شده است؛ یعنی گراف‌هایی که گره‌ها و یال‌های آن‌ها ویژگی متنی دارند \cite{he2024gretriever}. فرایند این روش چهار مرحله دارد. ابتدا متن گره‌ها و یال‌ها رمزگذاری و برای جست‌وجوی نزدیک‌ترین همسایه نمایه‌سازی می‌شود. سپس پرسش با همان رمزگذار به یک بردار تبدیل می‌گردد و مرتبط‌ترین گره‌ها و یال‌ها بر اساس شباهت با پرسش انتخاب می‌شوند. این امتیازها مشخص می‌کنند کدام اجزای گراف برای پاسخ اهمیت بیشتری دارند.

پس از امتیازدهی اولیه، انتخاب نهایی زیرگراف به صورت مسئله درخت اشتاینر جایزه‌جمع‌کن\footnote{\lr{Prize-Collecting Steiner Tree (PCST)}} صورت‌بندی می‌شود. در این صورت‌بندی، گره‌ها و یال‌های مرتبط جایزه بیشتری دارند و افزودن اجزای دیگر به زیرگراف با هزینه همراه است. نتیجه، زیرگرافی متصل است که میان پوشش اجزای مرتبط و محدود نگه‌داشتن اندازه ورودی تعادل برقرار می‌کند. زیرگراف به متن تبدیل و همراه با پرسش و نمایش حاصل از رمزگذار گراف به مدل زبانی داده می‌شود. بازگرداندن همین زیرگراف، امکان مشاهده شواهد انتخاب‌شده را نیز فراهم می‌کند.

مزیت ساختار \lr{G-Retriever} آن است که می‌تواند روی گرافی بزرگ‌تر از ظرفیت ورودی مدل عمل کند و تنها بخش مرتبط و متصل آن را وارد مرحله تولید نماید. در مقابل، استفاده از آن مستلزم وجود یک گراف متنی و محاسبه امتیاز مناسب برای گره‌ها و یال‌هاست. همچنین، کیفیت پاسخ به زیرگراف استخراج‌شده وابسته است؛ اگر جزء لازم در مرحله بازیابی امتیاز کافی نگیرد، در ورودی مدل نیز حضور نخواهد داشت. بنابراین، در بازپیاده‌سازی این روش باید عملکرد بازیاب و مولد به صورت جداگانه ثبت شود.

روش \lr{HippoRAG 2} با فرض متفاوتی آغاز می‌کند و گراف را از مجموعه اسناد می‌سازد \cite{gutierrez2025hipporag2}. در مرحله نمایه‌سازی برون‌خط، یک مدل زبانی با استخراج اطلاعات باز\footnote{\lr{Open Information Extraction (OpenIE)}} از هر قطعه، سه‌تایی‌های بدون شِمای ثابت تولید می‌کند. موضوع و مفعول این سه‌تایی‌ها به صورت گره‌های عبارتی و رابطه میان آن‌ها به صورت یال ذخیره می‌شود. سپس پیوندهای هم‌معنا بر پایه شباهت برداری افزوده می‌شوند. نسخه دوم \lr{HippoRAG} علاوه بر این مفاهیم، خود قطعه‌های متن را نیز به عنوان گره نگه می‌دارد و آن‌ها را به عبارت‌های استخراج‌شده متصل می‌کند؛ در نتیجه، ساختار مفهومی و زمینه اصلی متن در یک نمایه مشترک قرار می‌گیرند.

در زمان بازیابی، پرسش با سه‌تایی‌ها و قطعه‌های متن مقایسه می‌شود و نامزدهای اولیه برای شروع جست‌وجوی گرافی شکل می‌گیرند. یک مرحله بازشناسی با استفاده از مدل زبانی، سه‌تایی‌های نامرتبط را از میان نامزدهای برتر حذف می‌کند. سپس الگوریتم رتبه صفحه شخصی‌سازی‌شده\footnote{\lr{Personalized PageRank (PPR)}} با وزن‌دهی اولیه به گره‌های عبارتی و متنی اجرا می‌شود و قطعه‌ها بر اساس امتیاز نهایی رتبه‌بندی می‌گردند. خروجی بازیاب در این روش خود گراف نیست، بلکه قطعه‌های متنی برتری است که برای پاسخ‌گویی به مدل زبانی داده می‌شوند.

این طراحی، بازیابی واقعیت‌های مستقیم و پیوندهای چندگامی میان قطعه‌ها را در یک فرایند واحد دنبال می‌کند. با این حال، بخشی از هزینه در زمان ساخت گراف و اجرای مدل زبانی برای استخراج سه‌تایی‌ها صرف می‌شود و مرحله پالایش برخط نیز یک فراخوانی اضافی ایجاد می‌کند. در نتیجه، هنگام ارزیابی \lr{HippoRAG 2} باید هزینه آماده‌سازی نمایه از هزینه هر پرسش جدا شود. کیفیت سه‌تایی‌های استخراج‌شده و نحوه تنظیم وزن آغازین گره‌ها نیز بر رتبه نهایی قطعه‌ها اثر می‌گذارند.

روش \lr{LinearRAG} مسیر سبک‌تری برای ساخت گراف از متن در نظر می‌گیرد و به جای استخراج رابطه، سه نوع گره موجودیت، جمله و قطعه متن را در ساختاری به نام \lr{Tri-Graph} به هم متصل می‌کند \cite{zhuang2026linearrag}. بازیابی آن در دو مرحله انجام می‌شود: ابتدا موجودیت‌های مرتبط از طریق پیوندهای معنایی فعال می‌شوند و سپس اهمیت قطعه‌ها با انتشار امتیاز روی گراف محاسبه می‌گردد. این روش از جزئیات رابطه صرف‌نظر می‌کند تا ساخت نمایه ساده‌تر و کم‌هزینه‌تر باشد. با توجه به دامنه این پایان‌نامه، \lr{LinearRAG} بیشتر به عنوان نمونه‌ای از گراف‌سازی بدون رابطه صریح در نظر گرفته می‌شود و بررسی اصلی بر دو روش پیشین متمرکز خواهد بود.

در مجموع، \lr{G-Retriever} گراف را ورودی مسئله فرض می‌کند و یک زیرگراف متصل تحویل می‌دهد، در حالی که \lr{HippoRAG 2} و \lr{LinearRAG} ابتدا از پیکره متنی ساختار گرافی می‌سازند و در پایان قطعه‌های متن را بازیابی می‌کنند. تفاوت دیگر در سازوکار انتخاب است: اولی از بهینه‌سازی اتصال زیرگراف استفاده می‌کند و دو روش دیگر امتیاز را روی گراف منتشر می‌سازند. این تفاوت‌ها بر نوع داده مورد نیاز، هزینه نمایه‌سازی، امکان مشاهده مسیر بازیابی و حجم اطلاعات ورودی مدل اثر دارند. بخش بعد، خانواده عامل‌محور را بررسی می‌کند که در آن مدل زبانی در انتخاب گام‌های پیمایش نقش مستقیم‌تری دارد.

\section{روش‌های گرافی عامل‌محور: سازوکار و ویژگی‌ها}

در روش‌های عامل‌محور، مدل زبانی در یک چرخه جست‌وجو و ارزیابی شواهد قرار می‌گیرد. پرس‌وجو از گراف، همسایه‌ها و روابط موجود را مشخص می‌کند و مدل از میان این گزینه‌ها، ادامه‌های مرتبط با پرسش را برمی‌گزیند. پس از هر مرحله، اطلاعات گردآوری‌شده برای پاسخ‌گویی بررسی می‌شوند و در صورت ناکافی بودن، کاوش ادامه می‌یابد. دو روش \lr{ToG} و \lr{ToG-2} این الگو را به‌ترتیب با تکیه بر مسیرهای گرافی و با ترکیب گراف و متن به کار می‌گیرند.

روش \lr{ToG} جست‌وجو را از موجودیت‌های شناسایی‌شده در پرسش آغاز می‌کند و چند مسیر محتمل را به‌طور هم‌زمان نگه می‌دارد \cite{sun2024think}. این فرایند بر جست‌وجوی پرتوی\footnote{\lr{Beam Search}} استوار است؛ یعنی در هر مرحله، تنها تعداد محدودی از مسیرهای برتر برای گسترش بعدی حفظ می‌شوند. دو پارامتر اصلی، عرض جست‌وجو و بیشینه عمق هستند: اولی تعداد مسیرهای نگه‌داشته‌شده و دومی حداکثر تعداد گام‌های پیمایش را تعیین می‌کند. این محدودیت‌ها اندازه فضای جست‌وجو و هزینه تعامل با مدل را کنترل می‌کنند.

کاوش در \lr{ToG} دو مرحله دارد. ابتدا روابط متصل به انتهای هر مسیر از گراف خوانده می‌شوند و مدل، روابط مرتبط با پرسش را امتیازدهی و انتخاب می‌کند. سپس موجودیت‌های متصل از طریق روابط منتخب بازیابی می‌شوند و مدل از میان آن‌ها نامزدهای مناسب را برمی‌گزیند. بدین‌ترتیب، مسیرهای باقی‌مانده با یک سه‌تایی جدید گسترش می‌یابند. انتخاب رابطه پیش از موجودیت، تعداد نامزدهایی را که مدل باید بررسی کند کاهش می‌دهد؛ جهت یال نیز در بازیابی همسایه‌ها لحاظ می‌شود.

پس از گسترش مسیرها، مدل بررسی می‌کند که آیا اطلاعات موجود برای پاسخ کافی است. در صورت کفایت، پاسخ با استفاده از پرسش و مسیرها تولید می‌شود و در غیر این صورت، چرخه تا رسیدن به حد عمق ادامه می‌یابد. در نسخه اصلی، اگر جست‌وجو در این حد نیز شواهد کافی نیابد، پاسخ با اتکا به دانش درونی مدل تولید می‌شود. این حالت باید هنگام ارزیابی از پاسخ‌های متکی بر مسیر تفکیک شود. مسیرهای صریح، امکان مشاهده و بازبینی روابط پیموده‌شده را فراهم می‌کنند؛ اجرای روش نیز به تنظیم مجدد پارامترهای مدل زبانی نیاز ندارد.

روش \lr{ToG-2} همین جهت‌گیری تعاملی را با اسناد متنی متصل به موجودیت‌ها گسترش می‌دهد \cite{ma2025tog2}. پس از پیونددهی موجودیت‌های پرسش و انتخاب نقطه‌های شروع، ابتدا قطعه‌هایی از اسناد مرتبط بازیابی می‌شوند. اگر همین اطلاعات برای پاسخ کافی باشند، نیازی به پیمایش بیشتر نیست. در غیر این صورت، جست‌وجوی روابط و موجودیت‌های همسایه آغاز می‌شود تا منابع متنی تازه‌ای در دسترس قرار گیرند. در این طراحی، گراف محدوده بازیابی را هدایت می‌کند و متن، جزئیاتی را فراهم می‌سازد که لزوماً در سه‌تایی‌ها ثبت نشده‌اند.

تعامل متن و گراف در \lr{ToG-2} دوطرفه است. قطعه‌های مربوط به موجودیت‌های نامزد با در نظر گرفتن پرسش و رابطه‌ای که به آن‌ها منتهی شده رتبه‌بندی می‌شوند. سپس امتیاز قطعه‌ها برای انتخاب موجودیت‌های مرحله بعد به کار می‌رود؛ بنابراین، ادامه پیمایش فقط بر نام موجودیت یا رابطه متکی نیست. در پایان هر دور، مسیرها، قطعه‌های برتر و سرنخ‌های مراحل پیشین به مدل داده می‌شوند. مدل یا پاسخ را تولید می‌کند، یا سرنخ‌های مفید را خلاصه می‌سازد و با بازتنظیم پرسش، جست‌وجو را تا حد تعیین‌شده ادامه می‌دهد.

در مجموع، \lr{ToG} ساختاری ساده‌تر برای کاوش مسیرهای دانشی فراهم می‌کند و \lr{ToG-2} از متن برای تکمیل شواهد و هدایت دقیق‌تر کاوش بهره می‌گیرد. استفاده از نسخه دوم مستلزم ارتباط مشخص میان گراف و اسناد است و کیفیت این پیوندها بر بازیابی اثر دارد. در هر دو روش، افزایش عرض یا عمق می‌تواند شواهد بیشتری در دسترس قرار دهد، اما هزینه و زمان پاسخ را نیز تغییر می‌دهد. برای مقایسه در این پایان‌نامه، مسیرهای پیموده‌شده، شرط توقف، تعداد فراخوانی‌ها و شواهد نهایی ثبت خواهند شد تا کیفیت پاسخ در کنار نحوه دستیابی به آن بررسی شود.

\section{مقایسه روش‌های بازیابی‌محور و عامل‌محور}

تفاوت اصلی این دو خانواده در نحوه هدایت جست‌وجو است. در روش‌های بازیابی‌محور، الگوریتم بازیابی بر اساس پرسش، شواهد را انتخاب می‌کند و سپس آن‌ها را به مولد می‌سپارد. در روش‌های عامل‌محور، مدل زبانی با مشاهده شواهد هر مرحله درباره ادامه پیمایش و کافی بودن اطلاعات تصمیم می‌گیرد. بنابراین، صرف استفاده از مدل زبانی در ساخت گراف یا پالایش نامزدها، یک روش را عامل‌محور نمی‌کند؛ برای نمونه، مرحله بازشناسی در \lr{HippoRAG 2} بخشی از فرایند بازیابی است، اما در \lr{ToG-2} تصمیم‌های مدل بر ادامه کاوش گراف و متن اثر می‌گذارند \cite{gutierrez2025hipporag2,ma2025tog2}.

از نظر پوشش شواهد، بازیابی‌محور بودن به معنای محدود شدن به روابط مستقیم نیست. انتشار امتیاز در گراف یا استخراج یک زیرگراف متصل می‌تواند اطلاعات چندگامی را نیز گردآوری کند. مزیت روش عامل‌محور، امکان هدایت مرحله بعد با توجه به یافته‌های مرحله قبل است؛ این ویژگی زمانی مفید است که مفهوم واسط از ابتدا آشکار نباشد. در مقابل، انتخاب نامناسب یک رابطه یا حذف زودهنگام یک موجودیت می‌تواند ادامه جست‌وجو را محدود کند. در هر دو خانواده، پوشش گراف و کیفیت شناسایی مفاهیم اولیه بر شواهد قابل دسترس اثر دارند.

در توضیح‌پذیری نیز باید میان «شواهد انتخاب‌شده»، «مسیر پیمایش» و «توضیح پاسخ» تمایز گذاشت. زیرگراف خروجی \lr{G-Retriever} نشان می‌دهد چه اطلاعاتی در اختیار مدل بوده است و مسیرهای \lr{ToG} ترتیب روابط دنبال‌شده را آشکار می‌کنند \cite{he2024gretriever,sun2024think}. این آثار ثبت‌شده بازبینی پاسخ را آسان‌تر می‌کنند، اما به‌تنهایی نشان نمی‌دهند که تمام ادعاهای توضیح از همان شواهد نتیجه شده‌اند. همچنین، وجود یک مسیر در گراف با مرتبط بودن آن به پرسش یکسان نیست. از این رو، ردیابی‌پذیری جست‌وجو و وفاداری توضیح باید جداگانه ارزیابی شوند.

هزینه اجرا نیز به پیاده‌سازی هر روش وابسته است. در روش‌های بازیابی‌محور، بخشی از هزینه به ساخت نمایه یا آموزش مؤلفه گراف اختصاص دارد؛ روش‌های عامل‌محور معمولاً با افزایش گام‌های کاوش، فراخوانی‌های بیشتری در زمان پاسخ‌گویی دارند. تعداد فراخوانی به‌تنهایی معیار کافی نیست و طول ورودی و خروجی، اندازه مدل و زمان پردازش گراف نیز باید ثبت شوند. در مقایسه این پژوهش، داده و منابع دانش مشترک خواهند بود و تفاوت‌های لازم در نمایش گراف، آموزش و بودجه جست‌وجو گزارش می‌شوند. نتیجه مورد انتظار، مشخص شدن تناسب هر روش با نوع وظیفه و منابع در دسترس است؛ رتبه‌بندی نهایی به نتایج آزمایش‌ها وابسته خواهد بود.

\section{استدلال مقید به گراف: روش \lr{GCR}}

روش \lr{GCR} قید ساختاری را در زمان تولید مسیر اعمال می‌کند: مدل فقط می‌تواند دنباله‌ای بسازد که با یکی از مسیرهای مجاز گراف سازگار باشد \cite{luo2025gcr}. این روش، الگویی مستقل برای استدلال روی گراف دانش است و نباید آن را نسخه بعدی \lr{ToG-2} تلقی کرد. ارتباط آن با بحث پیشین در این است که انتخاب مسیر و تولید نمایش زبانی آن را به یکدیگر پیوند می‌دهد. چارچوب اصلی سه مرحله دارد: ساخت نمایه مسیرها، تولید مقید مسیر و پاسخ‌های نامزد، و ترکیب آن‌ها برای رسیدن به پاسخ نهایی.

در مرحله نخست، موجودیت‌های مطرح‌شده در پرسش به گراف پیوند داده می‌شوند و مسیرهای آغازشده از آن‌ها تا عمقی مشخص استخراج می‌گردند. مقاله برای این کار از جست‌وجوی سطح‌اول استفاده می‌کند. مسیرها پس از تبدیل به رشته‌های قالب‌بندی‌شده و نشانه‌های زبانی، در یک درخت پیشوندی\footnote{\lr{Trie}} با نام \lr{KG-Trie} ذخیره می‌شوند. پیشوندهای مشترک تنها یک بار در این ساختار قرار می‌گیرند و برای هر دنباله ناقص، ادامه‌های مجاز قابل بازیابی‌اند. بنابراین، نمایه شامل مسیرهای انتخاب‌شده در محدوده جست‌وجو است و لازم نیست همه مسیرهای ممکن گراف را در بر گیرد.

در مرحله رمزگشایی مقید\footnote{\lr{Constrained Decoding}}، مدل زبانی احتمال نشانه بعدی را محاسبه می‌کند، اما تنها نشانه‌هایی امکان انتخاب دارند که ادامه‌ای معتبر در \lr{KG-Trie} بسازند. این قید در طول تولید مسیر حفظ می‌شود تا رابطه یا موجودیتی خارج از مسیرهای نمایه‌شده وارد آن نشود. پس از پایان مسیر، رمزگشایی عادی برای تولید یک پاسخ نامزد از سر گرفته می‌شود. در نتیجه، دامنه قید ساختاری مشخص است: قید مستقیماً بر مسیر اعمال می‌شود و تمام متن پاسخ را در بر نمی‌گیرد.

در طراحی مقاله، یک مدل زبانی سبک‌تر با نمونه‌های پرسش، مسیر و پاسخ برای این مرحله تنظیم می‌شود. سپس چند مسیر و پاسخ نامزد با جست‌وجوی پرتوی تولید و به یک مدل زبانی عمومی داده می‌شوند تا پاسخ نهایی را با استفاده از آن‌ها بسازد. مدل تخصصی، انتخاب مسیرهای مرتبط را می‌آموزد و مدل عمومی ترکیب شواهد را انجام می‌دهد. این تفکیک نیازمند دسترسی به فرایند رمزگشایی مدل نخست است؛ از این رو، امکانات اجرایی آن با روشی که صرفاً از طریق ارسال درخواست متنی به مدل کار می‌کند متفاوت است.

تعبیر «نبود توهم» در این مقاله به انطباق کامل مسیر تولیدشده با گراف اشاره دارد. این ویژگی برای کنترل روابط ساختگی مفید است، اما صحت خود گراف، ارتباط مسیر با پرسش یا درستی پاسخ نهایی را تضمین نمی‌کند. گراف ناقص، پیونددهی نادرست موجودیت‌ها و محدودیت عمق جست‌وجو همچنان می‌توانند شواهد لازم را از دسترس خارج کنند. افزون بر این، توضیح نهایی ممکن است برداشتی فراتر از مسیرها داشته باشد. بنابراین، در انطباق با سلامت روان، اعتبار ساختاری مسیر و پشتیبانی ادعاهای توضیح از متن ورودی دو موضوع مجزا خواهند بود.

آزمایش‌های مقاله عمدتاً بر پرسش و پاسخ گراف دانش انجام شده‌اند و قابلیت انتقال به گراف‌های دیگر نیز بررسی شده است. این نتایج مبنایی برای انتخاب \lr{GCR} در مقایسه حاضر فراهم می‌کنند، اما عملکرد آن روی وظایف \lr{IMHI} باید مستقلاً سنجیده شود. هزینه ساخت \lr{KG-Trie}، تنظیم مدل تخصصی و اجرای دو مدل نیز در تحلیل کارایی منظور خواهد شد تا کاهش احتمالی رفت‌وبرگشت‌های جست‌وجو در کنار هزینه آماده‌سازی و تولید پاسخ بررسی شود.

\section{کاربرد در سلامت روان و مبنای ارزیابی}

کاربرد روش‌های مرورشده در این پژوهش، تولید برچسب و توضیح برای متن‌های سلامت روان است. در این وظیفه، پاسخ معمولاً یک موجودیت انتهایی در گراف نیست و باید با مجموعه برچسب‌های هر داده هماهنگ شود. مجموعه \lr{IMHI} چنین مسئله‌ای را در قالب دستور و پاسخ سازمان‌دهی می‌کند و مبنای آموزش و ارزیابی مدل‌های \lr{MentalLLaMA} بوده است \cite{yang2024mentallama}. وظایف منتخب شامل تشخیص استرس در \lr{Dreaddit}، طبقه‌بندی وضعیت در \lr{T-SID}، شناسایی عامل استرس در \lr{SAD} و بررسی ابعاد بهزیستی در \lr{MultiWD} هستند. این تنوع، امکان مقایسه روش‌ها در طبقه‌بندی دودویی و چندکلاسه، شناسایی عامل و تشخیص ابعاد بهزیستی را فراهم می‌کند.

در انتقال روش‌ها به این وظایف، متن فرد و دانش دامنه‌ای نقش‌های متفاوتی دارند. متن مشخص می‌کند چه تجربه یا نشانه‌ای گزارش شده است و گراف، روابط دانشی مربوط به مفاهیم آن را فراهم می‌کند. یک رابطه عمومی در گراف نمی‌تواند به‌تنهایی شاهدی بر وجود همان وضعیت در فرد باشد. برای مثال، اگر متن فقط به دشواری خواب اشاره کند، وجود روابط مربوط به خستگی در گراف مجوز نسبت دادن خستگی گزارش‌نشده به نویسنده نیست. بنابراین، حفظ بافت، نفی و میزان قطعیت عبارت‌ها در استخراج مفاهیم و تولید توضیح اهمیت دارد.

ارزیابی برچسب‌ها با دقت، بازخوانی\footnote{\lr{Recall}} و امتیاز \lr{F1} انجام خواهد شد. دقت نسبت نمونه‌های درست پیش‌بینی‌شده به همه نمونه‌های پیش‌بینی‌شده برای یک کلاس است و بازخوانی نسبت نمونه‌های درست شناسایی‌شده به همه نمونه‌های واقعی همان کلاس را نشان می‌دهد. امتیاز \lr{F1} میانگین هماهنگ دقت و بازخوانی است و زمانی بالا خواهد بود که هر دو معیار هم‌زمان مقدار مناسبی داشته باشند. برای یک کلاس، این معیارها به صورت زیر تعریف می‌شوند:
\begin{equation}
\mathrm{Precision}=\frac{TP}{TP+FP},\qquad
\mathrm{Recall}=\frac{TP}{TP+FN},\qquad
F1=2\times\frac{\mathrm{Precision}\times\mathrm{Recall}}
{\mathrm{Precision}+\mathrm{Recall}}.
\end{equation}
در این روابط، \lr{TP} نمونه‌های مثبت درست، \lr{FP} نمونه‌های منفی که به‌اشتباه مثبت تشخیص داده شده‌اند و \lr{FN} نمونه‌های مثبت ازدست‌رفته هستند. برای مجموعه‌های چندکلاسه، میانگین وزن‌دار هر معیار با وزن متناسب با تعداد نمونه‌های واقعی هر کلاس محاسبه می‌شود:
\begin{equation}
M_{\mathrm{weighted}}=\sum_{c=1}^{C}\frac{n_c}{N}M_c,
\end{equation}
که در آن \lr{M} می‌تواند دقت، بازخوانی یا \lr{F1} کلاس \lr{c} باشد؛ \lr{n} تعداد نمونه‌های آن کلاس و \lr{N} تعداد کل نمونه‌هاست. بنابراین، علاوه بر امتیازهای وزن‌دار، نتیجه هر کلاس نیز گزارش می‌شود تا ضعف عملکرد روی کلاس‌های کم‌نمونه پنهان نماند. این معیارها توافق خروجی با برچسب‌های مجموعه‌داده را می‌سنجند و نباید معادل اعتبار تشخیص بالینی تفسیر شوند.

برای توضیح‌ها، کیفیت کلی متن و پشتیبانی ادعاها جداگانه بررسی می‌شوند. در امتداد طرح پروپوزال، \lr{BERTScore} برای سنجش شباهت معنایی توضیح تولیدشده به توضیح مرجع در نظر گرفته می‌شود؛ این معیار با مقایسه نمایش‌های برداری واژه‌ها، میزان هم‌پوشانی معنایی دو متن را برآورد می‌کند. توضیح‌های مرجع \lr{IMHI} با کمک مدل زبانی تولید و ارزیابی شده‌اند، بنابراین شباهت به آن‌ها به‌تنهایی شاهد صحت همه ادعاها نیست. تطابق با متن فرد و دانش پشتیبان باید مستقل از این امتیاز سنجیده شود.

چارچوب \lr{GraphEval} برای این بررسی، خروجی مدل را به سه‌تایی‌های ادعایی تبدیل می‌کند و سازگاری هر سه‌تایی را با متن پشتیبان، با استفاده از مدل استنتاج زبان طبیعی، می‌سنجد \cite{sansford2024grapheval}. سه‌تایی‌های ناسازگار نیز مشخص می‌شوند تا محل خطا قابل بازبینی باشد. در کاربرد حاضر، ادعاهای مربوط به فرد باید با متن ورودی و ادعاهای دانشی با منابع دامنه‌ای مربوط بررسی شوند. این سازوکار با کنترل وجود یال‌های مسیر در گراف تفاوت دارد و نتیجه آن به دقت استخراج ادعا و ارزیاب وابسته است. نمونه‌هایی از خطاها نیز برای بررسی نفی، ابهام و تعمیم فراتر از شواهد به صورت کیفی تحلیل خواهند شد.

محور سوم، زمان پاسخ از دریافت ورودی تا تولید برچسب و توضیح، و مجموع نشانه‌های ورودی و خروجی در تمام فراخوانی‌هاست. هزینه ساخت گراف، نمایه‌سازی و آموزش از هزینه هر نمونه جدا گزارش می‌شود. استفاده از تقسیم‌های داده یکسان، حفظ مرز آموزش و آزمون و مشخص کردن بودجه منابع هر روش، مبنای مقایسه خواهد بود. این چارچوب نشان می‌دهد هر روش با چه هزینه‌ای به چه کیفیتی در برچسب و توضیح می‌رسد.

مرور این فصل نشان داد که روش‌های منتخب، نقش‌های متفاوتی به گراف می‌دهند: انتخاب شواهد در روش‌های بازیابی‌محور، هدایت جست‌وجوی مرحله‌ای در روش‌های عامل‌محور و محدود کردن تولید مسیر در \lr{GCR}. این تفاوت‌ها مبنای بازپیاده‌سازی و ارزیابی مشترک در سلامت روان خواهند بود. هدف آزمایش‌ها، سنجش اثر هر سازوکار بر درستی برچسب، پشتیبانی توضیح و هزینه اجراست. فصل بعد، جزئیات داده‌ها، آماده‌سازی گراف و پروتکل اجرای این مقایسه را مشخص می‌کند.
